% Source: https://artofproblemsolving.com/community/c3h1886571_geometry
% Stable identity: aops-topic-1886571-q01
Let $H$ be the midpoint of $BC$. Since $AB=AC$, the altitude $AH$ bisects the apex angle, so $\angle BAH=22.5^\circ$ and $BH=5$.

\begin{asy}
usepackage("amsmath");
size(4.5cm);
defaultpen(fontsize(10pt));
real h=5*(sqrt(2)+1);
pair A=(0,h),B=(-5,0),C=(5,0),H=(0,0);
draw(A--B--C--cycle,linewidth(1));
draw(A--H,dashed);
draw((0.6,0)--(0.6,0.6)--(0,0.6));
draw(arc(A,1.6,247.5,292.5));
label("$45^\circ$",A+(0.9,-0.9),E,fontsize(9pt));
label("$A$",A,N);label("$B$",B,SW);label("$C$",C,SE);label("$H$",H,S);
label("$5$",(B+H)/2,S);label("$5$",(H+C)/2,S);
label("$h$",(0,h*0.4),E);
label("\mbox{}",(min(currentpicture,user=true)+max(currentpicture,user=true))/2,black); // keeps the text after the figure black when LaTeX typesets the labels inline
\end{asy}

With $\theta=45^\circ$,
\[
\tan22.5^\circ=\dfrac{1-\cos\theta}{\sin\theta}=\dfrac{1-\dfrac{\sqrt2}{2}}{\dfrac{\sqrt2}{2}}=\sqrt2-1 .
\]
So the height is $h=AH=\dfrac{BH}{\tan22.5^\circ}=\dfrac{5}{\sqrt2-1}=5\left(\sqrt2+1\right)$, and the area is
\[
\dfrac{1}{2}\cdot10\cdot5\left(\sqrt2+1\right)=\boxed{25\left(1+\sqrt2\right)}\approx60.36 .
\]
(Check with the law of cosines: $AB^2=\dfrac{100}{2-\sqrt2}=50\left(2+\sqrt2\right)$ and $\dfrac{1}{2}AB^2\sin45^\circ=25\left(\sqrt2+1\right)$.)
